\section{Exact interfaces: arithmetic tests, fluid profiles, and source corrections}

\begin{proposition}[NN-A011: the specified fluid and branch channels preserve the A constraint]
Retain the exact source maps in workbench24 and workbench29i.
Under A, every finite matrix $[W(h_i,h_j)]$ of admissible scalar
Weil tests $h_i\in C_c^\infty(\mathbb R)$ is positive semidefinite,
including every polynomial channel $h_i=L_{P_i}U_i$ constructed in
workbench24. The branch-trace form on its retained two-sheet module
is a different form and can be negative without contradicting this.

In the released-witness arithmetic realization of workbench29i, take
$H=\Xi$ and
\[
 \Psi_{\nu,j}(t,q,\zeta)
 =\Xi(S_j(X_{\nu,t}(q))+\zeta).
\]
For every preterminal time, and almost every terminal label in the
proved terminal construction, all its zeros in the variable $\zeta$
are real under A, with unchanged multiplicity. This holds together
with the source's proved divergence of the specified fixed-label
partial-derivative norm. In this particular map the derivative
singularity cannot yield an off-critical zeta zero.
\end{proposition}
\begin{proof}
For A's real ordinates the exact Weil form is
$W(f,g)=\sum_\rho m_\rho\overline{\widehat f(\gamma_\rho)}
\widehat g(\gamma_\rho)$, absolutely convergent on the specified
test class. Hence
$c^*[W(h_i,h_j)]c=\sum_\rho m_\rho|
\widehat{\sum c_ih_i}(\gamma_\rho)|^2\ge0$.
The polynomial channel preserves that test class and compact support,
as proved by its explicit differential-operator formula in workbench24.
By contrast its branch trace has fibre matrix
$\operatorname{diag}(2,-4\lambda)$ at a real $\lambda>0$; it is
not the componentwise Weil form or its scalar pullback.

The shifts $a=S_j(X_{\nu,t}(q))$ are real at all the stated parameters.
The translation $\zeta\mapsto a+\zeta$ has inverse $z\mapsto z-a$,
and transforms a germ $(z-z_0)^mA(z)$ to
$(\zeta-(z_0-a))^mA(a+\zeta)$. It preserves multiplicity and
imaginary part exactly. A makes the zeros $z_0$ of $\Xi$ real.
The completed arithmetic argument is $s=1/2+i(a+\zeta)$,
so $\Re s=1/2-\Im\zeta=1/2$ at those zeros.
No estimate on $\partial_ta$ enters this identity, so its possible
divergence does not affect the conclusion.
\end{proof}

This closes the particular proposed counterexample route after A; it does
not assert a theorem about all hypothetical fluid-to-arithmetic maps.
The underlying categories are a concrete analytic test-function module,
its Hermitian Gram form, and translated analytic divisors. The branch
trace is an algebra trace in a quadratic algebra. The equality of those
forms was never available. Neither a PDE singularity nor this translation
creates a finite-dimensional pure motive, a negative Weil matrix, or a
complexity-class separation.

The Bost--Connes formulas in source22 also retain their domains. Its
physical Gibbs partition function is $\zeta(\beta)>0$ for real
$\beta>1$. A locates the nontrivial zeros of its meromorphic continuation
on $\Re s=1/2$. The source's strong-limit projection $P_1$ has continued
expectation $1/\zeta(s)$ and therefore poles at those zeros; finite
Euler-projector expectations are entire finite products and the strong
operator limit is not automatically a meromorphic limit through those
poles. A supplies zero location, not a positive Gibbs state at complex
temperature. Any ``Lee--Yang'' terminology must refer to this stated
analytic continuation and coordinate, not silently to a unit-circle
theorem for a finite statistical-mechanical partition polynomial.

\begin{proposition}[NN-A012: two exact source corrections in paper16]
The finite-gerbe path integral printed in
\texttt{thm:finite-gerbe-anomaly} satisfies
$\mathcal A_n(S^4)=n$. Consequently its later assertion
$\mathcal A(q_n)\otimes\mathcal A(-q_n)\simeq\mathbf1$
in \texttt{prop:anomaly-cancellation} fails for every $n>1$ with
that normalization. Complex conjugation is not the inverse of an
arbitrary non-unit-modulus partition value.

Also, the measure algebra of an arbitrary proper adelic curve is a
Boolean sigma-algebra, but need not be a complete Boolean algebra.
The completeness assertion in \texttt{thm:final-synthesis-v11}(3)
requires an additional measure-space hypothesis such as localizability.
\end{proposition}
\begin{proof}
For $S^4$ and $\Phi_n=\mathbb Z/n\mathbb Z$, one has
$|H^0|=n$, $|H^1|=1$, and $H^2=0$. The sole quadratic action is zero,
so the displayed source formula gives $n$ for both $q_n$ and $-q_n$.
Their product is $n^2$, whereas the tensor unit has value one. On this
example a compensating Euler factor $n^{-\chi(S^4)/2}=n^{-1}$
normalizes each value; a general inverse-theory assertion must specify
the full Euler normalization, not merely replace $q$ by $-q$.

For the second assertion, let $\Omega$ be uncountable and equip it
with the countable-or-cocountable sigma-algebra and counting measure.
Assign the trivial absolute value of any field to every point.
Every principal logarithm is identically zero, so this is a proper
adelic curve under the source's stated definition. The only null set
is empty. The displayed definition at the start of
\texttt{sec:adelic-support-defect} imposes integrability of principal
logarithms and the product formula, but no sigma-finiteness or
localizability; the counterexample is directed at that exact definition.
Thus the measure algebra
is exactly that sigma-algebra.
Choose an uncountable subset $A$ with uncountable complement. The
family of measurable singleton sets $\{\{a\}:a\in A\}$ has no
least measurable upper bound: any upper bound is cocountable and
contains some $x\notin A$, and deleting $x$ gives a smaller measurable
upper bound. Thus arbitrary suprema need not exist, whereas countable
suprema always do.
\end{proof}

These are finite or explicit set-theoretic counterexamples, not new
consequences of E or A. They are retained as source repairs because
otherwise the proposed geometry would rely on false normalization or
completeness assertions. They supply no computational lower bound.
